\phantomsection\label{sec:teaching-experience}
\section*{教学经历}
我的教学工作包括课程教学和学生指导，以基于问题的学习为指导原则。

\subsection*{课程教学}
我的课程教学遵循三项核心原则：扎实掌握所授内容，了解学生已有的经验和背景，以及贯彻基于问题的学习方法。
具体做法包括：编写贴合学生水平的博客式课程材料，开展富有启发性、连接理论与实际问题的课堂教学，以及循序渐进地安排实践练习。

\pubentry{奥尔堡大学, 2025 年至 2026 年}{\href{https://blog.yanlincs.com/ai-system/}{课程材料} \,/\, \href{https://github.com/orgs/AI-Systems-Infrastructure/repositories}{幻灯片}}
{AI Systems \& Infrastructure}
{}
{本科基础课程。讲授与人工智能模型的标准化交互，以及在实际硬件设施上部署可扩展的人工智能系统。}

\subsection*{学生指导}
我正持续积累指导经验，目前负责本科生和硕士生的学期项目。
指导重点：明确问题定义，批判性地分析设计选择，以及形成有说服力的结论。

\pubentry{奥尔堡大学, 2026}{\href{https://github.com/Logan-Lin/traj-rep-generalization/releases/latest/download/main.pdf}{项目提案}}
{Trajectory Representation Learning that Generalizes Across Regions, Modes, and Tasks}
{}
{硕士生学期项目。研究可跨下游任务、地理区域和交通方式泛化的轨迹表示学习方法，在道路网络和海底等深线等空间环境基础上对车辆和船舶轨迹进行建模，并处理不同交通方式之间差异巨大的采样间隔。}

\pubentry{奥尔堡大学, 2026}{\href{https://github.com/Logan-Lin/assess-data-traj/releases/latest/download/main.pdf}{项目提案}}
{Assessment of the Data Foundation for Trajectory Machine Learning}
{}
{硕士生学期项目。从以数据为中心的视角研究轨迹机器学习，评估训练数据质量，并开发数据过滤和增强方法，以提升模型在下游任务中的性能。}

\pubentry{奥尔堡大学, 2025 年至 2026 年}{\href{https://github.com/Logan-Lin/camera-sees-behind/releases/latest/download/proposal.pdf}{项目提案}}
{A Camera That Sees Behind}
{}
{本科生学期项目。开发可部署在嵌入式设备上的轻量模型，用于从图像或视频流中移除人物，构建符合 GDPR 要求的监控系统原型。}
